% ===== Джерела до розділу 4. Ключі з префіксом ch04: =====
% Позиції ядра (bib-core.tex) тут НЕ дублюються.

% --- Моделювання загроз ---
\bibitem{ch04:shostack} Shostack A. Threat Modeling: Designing for Security. Indianapolis: Wiley, 2014. 624~p.

\bibitem{ch04:schneier} Schneier B. Attack Trees. \emph{Dr. Dobb's Journal}. 1999. Vol.~24, No.~12. P.~21--29.

\bibitem{ch04:kordy} Kordy B., Piètre-Cambacédès L., Schweitzer P. DAG-based Attack and Defense Modeling: Don't Miss the Forest for the Attack Trees. \emph{Computer Science Review}. 2014. Vol.~13--14. P.~1--38. DOI: 10.1016/j.cosrev.2014.07.001.

\bibitem{ch04:pols} Pols P. The Unified Kill Chain: Designing a Unified Kill Chain for analyzing, comparing and defending against cyber attacks: Master thesis. Cyber Security Academy, 2017.

\bibitem{ch04:strom} Strom B. E. et al. MITRE ATT\&CK: Design and Philosophy. MITRE Technical Report MTR190121. McLean: The MITRE Corporation, 2020.

% --- Оцінювання ризиків ---
\bibitem{ch04:nist30} NIST. Guide for Conducting Risk Assessments. NIST SP 800-30, Rev.~1. Gaithersburg, MD: NIST, 2012. DOI: 10.6028/NIST.SP.800-30r1.

\bibitem{ch04:fair} Freund J., Jones J. Measuring and Managing Information Risk: A FAIR Approach. Oxford: Butterworth-Heinemann, 2014. 408~p.

\bibitem{ch04:gordonloeb} Gordon L. A., Loeb M. P. The Economics of Information Security Investment. \emph{ACM Transactions on Information and System Security}. 2002. Vol.~5, No.~4. P.~438--457. DOI: 10.1145/581271.581274.

\bibitem{ch04:cvss4} FIRST. Common Vulnerability Scoring System Version 4.0: Specification Document. Forum of Incident Response and Security Teams, 2023. URL: \url{https://www.first.org/cvss/v4.0/}.

\bibitem{ch04:epss} FIRST. Exploit Prediction Scoring System (EPSS). Forum of Incident Response and Security Teams. URL: \url{https://www.first.org/epss/}.

% --- Поширення та епідемічні моделі ---
\bibitem{ch04:staniford} Staniford S., Paxson V., Weaver N. How to 0wn the Internet in Your Spare Time. \emph{Proceedings of the 11th USENIX Security Symposium}. USENIX Association, 2002. P.~149--167.

\bibitem{ch04:zou} Zou C. C., Towsley D., Gong W. Email Worm Modeling and Defense. \emph{Proceedings of the 13th International Conference on Computer Communications and Networks (ICCCN)}. IEEE, 2004. P.~409--414.

% --- Українські джерела ---
\bibitem{ch04:ilienko} Ільєнко А. В., Ільєнко С. С., Яковенко О. Л., Галич Є. О., Павленко В. Г. Застосування штучного інтелекту для прогнозування кіберзагроз. \emph{Кібербезпека: освіта, наука, техніка}. 2024. URL: \url{https://csecurity.kubg.edu.ua/index.php/journal/article/view/631}.

\bibitem{ch04:lungol} Лунгол О. М. Багаторівнева модель атрибуції кібератак. \emph{Кібербезпека: освіта, наука, техніка}. 2025. Спеціальний випуск. URL: \url{https://csecurity.kubg.edu.ua/index.php/journal/article/view/904}.
